\documentclass[11pt]{article}
\usepackage{amsmath, amssymb, geometry, amsthm, enumitem, xcolor}
\geometry{letterpaper, margin=1in}

\title{\textbf{Homework 2: Logic (Implications and Biconditionals)}}
\author{Tashfeen Omran}
\date{September 2026}

\begin{document}
\maketitle

\tableofcontents
\newpage

\section{Part 1: Comprehensive Introduction, Context, and Prerequisites}

\subsection{Core Concepts}
In the study of formal mathematical logic, we transition from analyzing individual, isolated statements to analyzing compound statements. This chapter focuses specifically on \textbf{Implications} and \textbf{Biconditionals}.
\begin{itemize}
    \item \textbf{Implication ($P \Rightarrow Q$)}: Also known as a conditional statement, read as ``If $P$, then $Q$.'' Here, $P$ is the \textit{hypothesis} (or premise) and $Q$ is the \textit{conclusion}. Crucially, an implication is only false when a true hypothesis leads to a false conclusion. If the hypothesis is false, the statement is considered \textit{vacuously true}, regardless of the truth value of $Q$.
    \item \textbf{Biconditional ($P \Leftrightarrow Q$)}: Read as ``$P$ if and only if $Q$'' (often abbreviated as ``iff''). This is the logical conjunction of two implications: $(P \Rightarrow Q) \wedge (Q \Rightarrow P)$. A biconditional is true if and only if $P$ and $Q$ share the exact same truth value (both true or both false).
    \item \textbf{Open Sentences over a Domain}: Sentences containing variables, such as $P(x): x^2 > 4$, are not propositions themselves until a specific value from a domain $S$ is substituted for $x$. 
\end{itemize}

\subsection{Intuition and Derivation}
The logic of $P \Rightarrow Q$ can be intuitively understood as a promise or a contract. Suppose a professor says, ``If you score a 100 on the final ($P$), then you will get an A in the class ($Q$).'' 
\begin{itemize}
    \item If you score 100 ($P$ is True) and get an A ($Q$ is True), the promise was kept. ($T \Rightarrow T$ is True).
    \item If you score 100 ($P$ is True) but get a B ($Q$ is False), the promise was broken. ($T \Rightarrow F$ is False).
    \item If you score an 80 ($P$ is False), the professor did not specify what would happen. Whether you get an A or a B, you cannot claim the professor lied. Thus, when $P$ is False, the conditional promise is not violated. ($F \Rightarrow T$ is True, and $F \Rightarrow F$ is True).
\end{itemize}

\subsection{Historical Context and Motivation}
The formalization of logical rules dates back to Aristotle's syllogisms, but the algebraic treatment of logic we use today was pioneered by George Boole in his 1854 work, \textit{The Laws of Thought}. Later mathematicians, such as Gottlob Frege and Bertrand Russell, extended these ideas to create rigorous systems capable of modeling all of mathematics. Mathematicians needed absolute precision in defining when a theorem ``If $A$, then $B$'' is true, leading to the strict truth-table definitions used today, completely devoid of subjective human interpretation.

\subsection{Key Formulas and Truth Tables}
The truth tables form the axioms of propositional logic:
\[
\begin{array}{|c|c|c|c|c|c|}
\hline
P & Q & \sim P & P \wedge Q & P \vee Q & P \Rightarrow Q & P \Leftrightarrow Q \\
\hline
T & T & F & T & T & T & T \\
T & F & F & F & T & F & F \\
F & T & T & F & T & T & F \\
F & F & T & F & F & T & T \\
\hline
\end{array}
\]
Other fundamental formulations:
\begin{align*}
    P \Rightarrow Q &\equiv (\sim P) \vee Q \\
    P \Leftrightarrow Q &\equiv (P \Rightarrow Q) \wedge (Q \Rightarrow P)
\end{align*}

\subsection{Prerequisites}
To master this topic, you must be comfortable with:
\begin{itemize}
    \item \textbf{Set Theory Notation}: Understanding domains, intersections ($\cap$), unions ($\cup$), elements ($\in, \notin$), and partitions.
    \item \textbf{Basic Algebra}: Solving linear/quadratic inequalities (e.g., $x^2 \geq 1 \Rightarrow x \in (-\infty, -1] \cup [1, \infty)$).
    \item \textbf{Number Theory Basics}: Recognizing prime numbers, even numbers (form $2k$), and odd numbers (form $2k+1$).
\end{itemize}

\newpage
\section{Part 2: Detailed Homework Solutions}

\subsection{Section 2.4: Implications}

\textbf{Exercise 2.19} \\
Consider the statements $P$: 17 is even, and $Q$: 19 is prime. Write each in words and indicate true or false.
\begin{itemize}
    \item We note that $P$ is False ($17$ is odd) and $Q$ is True ($19$ has no divisors other than $1$ and $19$).
    \item \textbf{(a) $\sim P$}: ``$17$ is not even.'' (True, since $P$ is False).
    \item \textbf{(b) $P \vee Q$}: ``$17$ is even or $19$ is prime.'' (True, since $Q$ is True).
    \item \textbf{(c) $P \wedge Q$}: ``$17$ is even and $19$ is prime.'' (False, since $P$ is False).
    \item \textbf{(d) $P \Rightarrow Q$}: ``If $17$ is even, then $19$ is prime.'' (True, because False $\Rightarrow$ True is True).
\end{itemize}

\vspace{10pt}
\textbf{Exercise 2.20} \\
Construct a truth table for $(P \Rightarrow Q) \Rightarrow (\sim P)$.
\[
\begin{array}{|c|c|c|c|c|}
\hline
P & Q & P \Rightarrow Q & \sim P & (P \Rightarrow Q) \Rightarrow (\sim P) \\
\hline
T & T & T & F & \textbf{F} \\
T & F & F & F & \textbf{T} \\
F & T & T & T & \textbf{T} \\
F & F & T & T & \textbf{T} \\
\hline
\end{array}
\]

\vspace{10pt}
\textbf{Exercise 2.21} \\
Let $P$: $\sqrt{2}$ is rational, and $Q$: $22/7$ is rational. We know $P$ is False and $Q$ is True.
\begin{itemize}
    \item \textbf{(a) $P \Rightarrow Q$}: False $\Rightarrow$ True. \textbf{Result: True}.
    \item \textbf{(b) $Q \Rightarrow P$}: True $\Rightarrow$ False. \textbf{Result: False}.
    \item \textbf{(c) $(\sim P) \Rightarrow (\sim Q)$}: True $\Rightarrow$ False. \textbf{Result: False}.
    \item \textbf{(d) $(\sim Q) \Rightarrow (\sim P)$}: False $\Rightarrow$ True. \textbf{Result: True}.
\end{itemize}

\vspace{10pt}
\textbf{Exercise 2.22} \\
Let $P$: $\sqrt{2}$ is rational (False), $Q$: $2/3$ is rational (True), $R$: $\sqrt{3}$ is rational (False).
\begin{itemize}
    \item \textbf{(a) $(P \wedge Q) \Rightarrow R$}: $(F \wedge T) \Rightarrow F \equiv F \Rightarrow F$. \textbf{Result: True}.
    \item \textbf{(b) $(P \wedge Q) \Rightarrow (\sim R)$}: $F \Rightarrow T$. \textbf{Result: True}.
    \item \textbf{(c) $((\sim P) \wedge Q) \Rightarrow R$}: $(T \wedge T) \Rightarrow F \equiv T \Rightarrow F$. \textbf{Result: False}.
    \item \textbf{(d) $(P \vee Q) \Rightarrow (\sim R)$}: $(F \vee T) \Rightarrow T \equiv T \Rightarrow T$. \textbf{Result: True}.
\end{itemize}

\vspace{10pt}
\textbf{Exercise 2.23} \\
Suppose $\{S_1, S_2\}$ is a partition of $S$ and $x \in S$. By definition of a partition, $S_1 \cap S_2 = \emptyset$ and $S_1 \cup S_2 = S$.
\begin{itemize}
    \item \textbf{(a) If $x \notin S_1$, then $x \in S_2$}: True. Since $x$ is in $S$ but not $S_1$, it must be in $S_2$.
    \item \textbf{(b) It's possible that $x \notin S_1$ and $x \notin S_2$}: False. The union is $S$.
    \item \textbf{(c) Either $x \notin S_1$ or $x \notin S_2$}: True. Since they are disjoint, $x$ cannot be in both, so it must be absent from at least one.
    \item \textbf{(d) Either $x \in S_1$ or $x \in S_2$}: True. Their union forms $S$.
    \item \textbf{(e) It's possible that $x \in S_1$ and $x \in S_2$}: False. The sets are disjoint.
\end{itemize}

\vspace{10pt}
\textbf{Exercise 2.24} \\
Two sets $A$ and $B$ are nonempty disjoint subsets of $S$. Let $x \in S$.
\begin{itemize}
    \item \textbf{(a) It's possible that $x \in A \cap B$}: False. They are disjoint ($A \cap B = \emptyset$).
    \item \textbf{(b) If $x \in A$, then $x \notin B$}: True. Because they are disjoint.
    \item \textbf{(c) If $x \notin A$, then $x \in B$}: False. They do not necessarily partition $S$.
    \item \textbf{(d) It's possible that $x \notin A$ and $x \notin B$}: True. $x$ could be in $S \setminus (A \cup B)$.
    \item \textbf{(e) For each nonempty set $C$, either $x \in A \cap C$ or $x \in B \cap C$}: False. $C$ might not contain $x$.
    \item \textbf{(f) For some nonempty set $C$, both $x \in A \cup C$ and $x \in B \cup C$}: True. Let $C = S$, then $A \cup C = S$ and $B \cup C = S$. Since $x \in S$, this holds.
\end{itemize}

\vspace{10pt}
\textbf{Exercise 2.25} \\
Statement: If I receive an A in both Calc I ($P$) and DM ($Q$), then I will take Calc II ($R$) or CP ($S$). Formally: $(P \wedge Q) \Rightarrow (R \vee S)$.
\begin{itemize}
    \item \textbf{(a)} $\sim P \wedge R$. The premise $P \wedge Q$ is False. Implication is True. (True)
    \item \textbf{(b)} $P \wedge Q \wedge \sim R \wedge \sim S$. True premise, False conclusion. (False)
    \item \textbf{(c)} $\sim P \wedge \sim R \wedge S$. Premise False. Implication is True. (True)
    \item \textbf{(d)} $P \wedge Q \wedge R \wedge S$. Premise True, Conclusion $(T \vee T)$ is True. (True)
    \item \textbf{(e)} $\sim P \wedge \sim Q \wedge \sim R \wedge \sim S$. Premise False. Implication is True. (True)
\end{itemize}

\vspace{10pt}
\textbf{Exercise 2.26} \\
Statement: If I don't see my advisor today ($\sim P$), then I'll see her tomorrow ($Q$). Formally: $\sim P \Rightarrow Q$.
\begin{itemize}
    \item \textbf{(a)} Doesn't see either day. $\sim P$ is True, $Q$ is False. True $\Rightarrow$ False. \textbf{False}.
    \item \textbf{(b)} Sees both days. $\sim P$ is False, $Q$ is True. False $\Rightarrow$ True. \textbf{True}.
    \item \textbf{(c)} Sees on one of the days. If today, $\sim P$ is False, result True. If tomorrow, $Q$ is True, result True. \textbf{True}.
    \item \textbf{(d)} Sees today ($P$ is True $\Rightarrow \sim P$ is False) and waits till next week ($Q$ is False). False $\Rightarrow$ False. \textbf{True}. (Wait, the question says ``Doesn't see his advisor today...'' Wait, option (d) says ``The student doesn't see his advisor today and waits until next week''. This means $\sim P$ is True, and $Q$ is False. True $\Rightarrow$ False. \textbf{False}.)
\end{itemize}

\vspace{10pt}
\textbf{Exercise 2.27} \\
Given conditions: Alice attends $\Rightarrow$ Ben attends. Ben attends $\Rightarrow$ Cindy attends. Cindy attends $\Rightarrow$ Don attends. Exactly two students attend.
Let the variables be $A, B, C, D$. If $A$ is True, then $B, C, D$ must be True, making 4 attendees (Contradicts ``exactly two''). Thus, $A$ is False. If $B$ is True, then $C, D$ are True, making 3 attendees. Thus, $B$ is False. For exactly two attendees, the remaining $C$ and $D$ must be True. Check: $C \Rightarrow D$ holds ($T \Rightarrow T$). The \textbf{two students are Cindy and Don}.

\vspace{10pt}
\textbf{Exercise 2.28} \\
Statement: $P \Rightarrow Q$. We want to find which option mathematically guarantees that $P \Rightarrow Q$ is True.
\begin{itemize}
    \item \textbf{(a) $Q \Rightarrow P$}: No.
    \item \textbf{(b) $\sim P \vee Q$}: Yes, this is logically equivalent to $P \Rightarrow Q$.
    \item \textbf{(c) $P$}: No, $Q$ could still be False.
    \item \textbf{(d) $P \wedge Q$}: Yes. Both are true, so $T \Rightarrow T$ is True.
    \item \textbf{(e) $P \wedge \sim Q$}: No, this makes it False.
    \item \textbf{(f) The concert is canceled}: Yes. If canceled, $P$ (Bill takes Sam to concert) is False. If $P$ is False, $P \Rightarrow Q$ is vacuously True.
    \item \textbf{(g) Sam doesn't attend}: Yes, same logic as (f), $P$ becomes False.
\end{itemize}

\vspace{10pt}
\textbf{Exercise 2.29} \\
Which implies $P \vee Q$ is False? (This requires both $P$ and $Q$ to be False).
\begin{itemize}
    \item \textbf{(a) $(\sim P) \vee (\sim Q)$ is False}: This implies $P$ and $Q$ are both True. No.
    \item \textbf{(b) $(\sim P) \vee Q$ is True}: $P=F, Q=T$ works here, but $P \vee Q$ would be True. No.
    \item \textbf{(c) $(\sim P) \wedge (\sim Q)$ is True}: This explicitly means $P$ is False and $Q$ is False. This completely implies $P \vee Q$ is False. \textbf{Yes.}
    \item \textbf{(d) $Q \Rightarrow P$ is True}: $Q=F, P=T$ works, making $P \vee Q$ True. No.
    \item \textbf{(e) $P \wedge Q$ is False}: $P=T, Q=F$ works, making $P \vee Q$ True. No.
\end{itemize}

\subsection{Section 2.5: More on Implications}

\textbf{Exercise 2.30} \\
$P(n): 5n+3$ is prime. $Q(n): 7n+1$ is prime. $n \in \mathbb{N}$.
\begin{itemize}
    \item \textbf{(a)} $P(n) \Rightarrow Q(n)$ in words: ``If $5n+3$ is prime, then $7n+1$ is prime.''
    \item \textbf{(b)} $n=2$: $P(2) = 13$ (Prime, True). $Q(2) = 15$ (Not prime, False). $T \Rightarrow F$ is \textbf{False}.
    \item \textbf{(c)} $n=6$: $P(6) = 33$ (Not prime, False). $Q(6) = 43$ (Prime, True). $F \Rightarrow T$ is \textbf{True}.
\end{itemize}

\vspace{10pt}
\textbf{Exercise 2.31} \\
\begin{itemize}
    \item \textbf{(a)} $P(x): |x|=4, Q(x): x=4, S=\{-4,-3,1,4,5\}$.
    For $x=-4$: $P(-4)$ is True, $Q(-4)$ is False. $T \Rightarrow F$ is False. Thus, the open sentence implication over $S$ is \textbf{False} for some $x$, specifically $x=-4$.
    \item \textbf{(b)} $P(x): x^2=16, Q(x): |x|=4, S=\{-6,-4,0,3,4,8\}$.
    Since $x^2=16 \Leftrightarrow |x|=4$ for all reals, $P(x)$ and $Q(x)$ always have the same truth value. Thus $P(x) \Rightarrow Q(x)$ is always \textbf{True}.
    \item \textbf{(c)} $P(x): x>3, Q(x): 4x-1 > 12 \Leftrightarrow x > 3.25, S=\{0,2,3,4,6\}$.
    For $x \in \{0,2,3\}$, $P(x)$ is False, making it vacuously True. For $x \in \{4,6\}$, $P(x)$ is True, and $Q(x)$ is also True. It evaluates to \textbf{True} for all $x \in S$.
\end{itemize}

\vspace{10pt}
\textbf{Exercise 2.32} \\
Determine all $x \in S$ for which $P(x) \Rightarrow Q(x)$ is True.
\begin{itemize}
    \item \textbf{(a)} $P(x): x=7, Q(x): x \geq 8, S=\mathbb{R}$.
    If $x=7$, $P$ is True, $Q$ is False (Implication False). For $x \neq 7$, $P$ is False (Implication True). Answer: \textbf{$S \setminus \{7\}$}.
    \item \textbf{(b)} $P(x): x^2 \geq 1, Q(x): x \geq 1, S=\mathbb{R}$.
    $P$ is True for $x \in (-\infty, -1] \cup [1, \infty)$. $Q$ is True for $x \in [1, \infty)$. The implication is False when $P$ is True and $Q$ is False, which occurs on $(-\infty, -1]$. The implication is True everywhere else. Answer: \textbf{$(-1, \infty)$}.
    \item \textbf{(c)} $P(x): x^2 \geq 1, Q(x): x \geq 1, S=\mathbb{N}$.
    Since $\mathbb{N} = \{1, 2, \dots\}$, $x \geq 1$ is always True. Both are always True. Answer: \textbf{$\mathbb{N}$}.
    \item \textbf{(d)} $P(x): x \in [-1, 2], Q(x): x^2 \leq 2, S=[-1, 1]$.
    For any $x \in S$, $P(x)$ is True, and $x^2 \leq 1 \leq 2$, so $Q(x)$ is also True. $T \Rightarrow T$. Answer: \textbf{$S = [-1, 1]$}.
\end{itemize}

\vspace{10pt}
\textbf{Exercise 2.33} \\
Evaluate $P(x,y) \Rightarrow Q(x,y)$ on given points.
\begin{itemize}
    \item \textbf{(a)} $P: x^2-y^2=0 \Leftrightarrow |x|=|y|$. $Q: x=y$.
    $(1,-1): P$ is True, $Q$ is False $\Rightarrow$ \textbf{False}.
    $(3,4): P$ is False $\Rightarrow$ \textbf{True}.
    $(5,5): P$ is True, $Q$ is True $\Rightarrow$ \textbf{True}.
    
    \item \textbf{(b)} $P: |x|=|y|, Q: x=y$.
    $(1,2): P$ is False $\Rightarrow$ \textbf{True}.
    $(2,-2): P$ is True, $Q$ is False $\Rightarrow$ \textbf{False}.
    $(6,6): P$ is True, $Q$ is True $\Rightarrow$ \textbf{True}.
    
    \item \textbf{(c)} $P: x^2+y^2=1, Q: x+y=1$.
    $(1,-1): P$ is False ($2 \neq 1$) $\Rightarrow$ \textbf{True}.
    $(-3,4): P$ is False ($25 \neq 1$) $\Rightarrow$ \textbf{True}.
    $(0,-1): P$ is True, $Q$ is False ($-1 \neq 1$) $\Rightarrow$ \textbf{False}.
    $(1,0): P$ is True, $Q$ is True $\Rightarrow$ \textbf{True}.
\end{itemize}

\vspace{10pt}
\textbf{Exercise 2.34} \\
Write as an ``if, then'' implication.
\begin{itemize}
    \item \textbf{(a)} If a point lies on the line $2y+x-3=0$ and its x-coordinate is an integer, then its y-coordinate is also an integer.
    \item \textbf{(b)} If an integer is odd, then its square is odd.
    \item \textbf{(c)} If $n \in \mathbb{Z}$ and $3n+7$ is even, then $n$ is odd.
    \item \textbf{(d)} If a function is $f(x) = \cos x$, then its derivative is $f'(x) = -\sin x$.
    \item \textbf{(e)} If $C$ is a circle with circumference $4\pi$, then the area of $C$ is $4\pi$.
    \item \textbf{(f)} If the integer $n^3$ is even, then $n$ is even. (Recall ``$P$ only if $Q$'' means $P \Rightarrow Q$).
\end{itemize}

\subsection{Section 2.6: Biconditionals}

\textbf{Exercise 2.35} \\
$P$: $18$ is odd (False), $Q$: $25$ is even (False). $P \Leftrightarrow Q$ is ``$18$ is odd if and only if $25$ is even.'' Since both are False, the biconditional is \textbf{True}.

\vspace{10pt}
\textbf{Exercise 2.36} \\
$P(x): x$ is odd, $Q(x): x^2$ is odd, over $\mathbb{Z}$.
\begin{itemize}
    \item \textbf{(1)} An integer $x$ is odd if and only if its square $x^2$ is odd.
    \item \textbf{(2)} An integer $x$ being odd is a necessary and sufficient condition for $x^2$ to be odd.
\end{itemize}

\vspace{10pt}
\textbf{Exercise 2.37} \\
$P(x): |x-3| < 1$, $Q(x): x \in (2,4)$ over $\mathbb{R}$.
\begin{itemize}
    \item \textbf{(1)} The inequality $|x-3| < 1$ holds if and only if $x \in (2, 4)$.
    \item \textbf{(2)} A necessary and sufficient condition for $|x-3| < 1$ is that $x$ lies in the interval $(2, 4)$.
\end{itemize}

\vspace{10pt}
\textbf{Exercise 2.38} \\
$S=\{-2,0,2\}, P(x): x=-2, Q(x): x^2=4$. Find $x$ for which statements are True.
\begin{itemize}
    \item \textbf{(a)} $\sim P(x)$: $x \neq -2$. True for \textbf{$\{0, 2\}$}.
    \item \textbf{(b)} $P(x) \vee Q(x)$: $x=-2$ or $x^2=4$. $Q$ covers both $2$ and $-2$. True for \textbf{$\{-2, 2\}$}.
    \item \textbf{(c)} $P(x) \wedge Q(x)$: Both must hold. True for \textbf{$\{-2\}$}.
    \item \textbf{(d)} $P(x) \Rightarrow Q(x)$: If $x=-2$, $x^2=4$ is True. If $x \neq -2$, $P$ is False (vacuously true). True for \textbf{$\{-2, 0, 2\}$}.
    \item \textbf{(e)} $Q(x) \Rightarrow P(x)$: Fails at $x=2$ (True $\Rightarrow$ False). True for \textbf{$\{-2, 0\}$}.
    \item \textbf{(f)} $P(x) \Leftrightarrow Q(x)$: Must share truth value. True for \textbf{$\{-2, 0\}$}.
\end{itemize}

\vspace{10pt}
\textbf{Exercise 2.39} \\
Determine $x \in S$ for which $P(x) \Leftrightarrow Q(x)$ is true.
\begin{itemize}
    \item \textbf{(a)} $P(x): |x|=4, Q(x): x=4, S=\{-4,-3,1,4,5\}$.
    Fails at $x=-4$ where $P$ is True but $Q$ is False. Matches everywhere else. Answer: \textbf{$\{-3, 1, 4, 5\}$}.
    \item \textbf{(b)} $P(x): x \geq 3, Q(x): 4x-1 > 12 \Leftrightarrow x > 3.25, S=\{0,2,3,4,6\}$.
    At $x=3$, $P$ is True, $Q$ is False. Matches at all other points. Answer: \textbf{$\{0, 2, 4, 6\}$}.
    \item \textbf{(c)} $P(x): x^2=16 \Leftrightarrow x \in \{-4,4\}, Q(x): x^2-4x=0 \Leftrightarrow x \in \{0,4\}, S=\{-6,-4,0,3,4,8\}$.
    At $x=-4$, $P$ is True, $Q$ is False. At $x=0$, $P$ is False, $Q$ is True. They match at $x=4$ (both True) and $x \in \{-6,3,8\}$ (both False). Answer: \textbf{$\{-6, 3, 4, 8\}$}.
\end{itemize}

\vspace{10pt}
\textbf{Exercise 2.40} \\
Truth value of $P(x,y) \Leftrightarrow Q(x,y)$.
\begin{itemize}
    \item \textbf{(a)} $P: x^2-y^2=0, Q: x=y$.
    $(1,-1)$: $P$ is True, $Q$ is False $\Rightarrow$ \textbf{False}.
    $(3,4)$: $P$ is False, $Q$ is False $\Rightarrow$ \textbf{True}.
    $(5,5)$: $P$ is True, $Q$ is True $\Rightarrow$ \textbf{True}.
    \item \textbf{(b)} $P: |x|=|y|, Q: x=y$.
    $(1,2)$: Both False $\Rightarrow$ \textbf{True}.
    $(2,-2)$: $P$ is True, $Q$ is False $\Rightarrow$ \textbf{False}.
    $(6,6)$: Both True $\Rightarrow$ \textbf{True}.
    \item \textbf{(c)} $P: x^2+y^2=1, Q: x+y=1$.
    $(1,-1)$: Both False $\Rightarrow$ \textbf{True}.
    $(-3,4)$: Both False $\Rightarrow$ \textbf{True}.
    $(0,-1)$: $P$ is True, $Q$ is False $\Rightarrow$ \textbf{False}.
    $(1,0)$: Both True $\Rightarrow$ \textbf{True}.
\end{itemize}

\vspace{10pt}
\textbf{Exercise 2.41} \\
$S=\{1,2,3\}$. $P(n): \frac{n^2+n}{2}$ is even. $Q(n): \frac{n^2-n}{2}$ is odd. We want $n$ where $P(n) \Leftrightarrow Q(n)$ is True.
\begin{itemize}
    \item $n=1: P(1) = 1$ (odd, False). $Q(1) = 0$ (even, False). $F \Leftrightarrow F$ is \textbf{True}.
    \item $n=2: P(2) = 3$ (odd, False). $Q(2) = 1$ (odd, True). $F \Leftrightarrow T$ is \textbf{False}.
    \item $n=3: P(3) = 6$ (even, True). $Q(3) = 3$ (odd, True). $T \Leftrightarrow T$ is \textbf{True}.
\end{itemize}
Answer: \textbf{$\{1, 3\}$}.

\vspace{10pt}
\textbf{Exercise 2.42} \\
$S=\{2,3,4\}$. $P(n): \frac{n(n-1)}{2}$ is odd. $Q(n): \frac{n(n+1)}{2}$ is even. We want $n$ where $P(n) \Leftrightarrow Q(n)$ is True.
\begin{itemize}
    \item $n=2: P(2) = 1$ (True). $Q(2) = 3$ (False). $T \Leftrightarrow F$ is \textbf{False}.
    \item $n=3: P(3) = 3$ (True). $Q(3) = 6$ (True). $T \Leftrightarrow T$ is \textbf{True}.
    \item $n=4: P(4) = 6$ (False). $Q(4) = 10$ (True). $F \Leftrightarrow T$ is \textbf{False}.
\end{itemize}
Answer: \textbf{$\{3\}$}.

\vspace{10pt}
\textbf{Exercise 2.43} \\
$S=\{1,2,3\}$. $P(n): \frac{(n+4)(n+5)}{2}$ is odd. $Q(n): 2^{n-2} + 3^{n-2} + 6^{n-2} > (2.5)^{n-1}$.
\begin{itemize}
    \item $n=1$: $P(1) = 15$ (True). $Q(1): 1/2 + 1/3 + 1/6 = 1$. $1 > 1$ is False.
    \item $n=2$: $P(2) = 21$ (True). $Q(2): 1 + 1 + 1 = 3$. $3 > 2.5$ is True.
    \item $n=3$: $P(3) = 28$ (False). $Q(3): 2 + 3 + 6 = 11$. $11 > 6.25$ is True.
\end{itemize}
We need:
\begin{itemize}
    \item $P(a) \Rightarrow Q(a)$ is False $\implies T \Rightarrow F$. Thus \textbf{$a=1$}.
    \item $Q(b) \Rightarrow P(b)$ is False $\implies T \Rightarrow F$. Thus \textbf{$b=3$}.
    \item $P(c) \Leftrightarrow Q(c)$ is True $\implies$ same truth value. Thus \textbf{$c=2$}.
\end{itemize}

\vspace{10pt}
\textbf{Exercise 2.44} \\
$S=\{1,2,3,4\}$. $P(n): \frac{n(n-1)}{2}$ is even. $Q(n): 2^{n-2} - (-2)^{n-2}$ is even. $R(n): 5^{n-1} + 2^n$ is prime.
\begin{itemize}
    \item $n=1: P(1)=0 \text{ (T)}, Q(1)=1 \text{ (F)}, R(1)=3 \text{ (T)}$.
    \item $n=2: P(2)=1 \text{ (F)}, Q(2)=0 \text{ (T)}, R(2)=9 \text{ (F)}$.
    \item $n=3: P(3)=3 \text{ (F)}, Q(3)=4 \text{ (T)}, R(3)=33 \text{ (F)}$.
    \item $n=4: P(4)=6 \text{ (T)}, Q(4)=0 \text{ (T)}, R(4)=141 \text{ (F, divisible by 3)}$.
\end{itemize}
Find distinct $a,b,c,d$:
\begin{itemize}
    \item $P(a) \Rightarrow Q(a)$ False $\implies a=1$.
    \item $Q(b) \Rightarrow P(b)$ True $\implies b \in \{1,4\}$. Since distinct, \textbf{$b=4$}.
    \item $P(c) \Leftrightarrow R(c)$ True $\implies c \in \{2,3\}$.
    \item $Q(d) \Leftrightarrow R(d)$ False $\implies d \in \{2,3\}$.
\end{itemize}
Valid distinct combinations are \textbf{$a=1, b=4, c=2, d=3$} or \textbf{$a=1, b=4, c=3, d=2$}.

\vspace{10pt}
\textbf{Exercise 2.45} \\
$S=\{2,3,4,5,6,11\}$. Statement: $2^n-1$ is prime $\Leftrightarrow n$ is prime.
\begin{itemize}
    \item $n=2$: $3$ is prime $\Leftrightarrow 2$ is prime. (True)
    \item $n=3$: $7$ is prime $\Leftrightarrow 3$ is prime. (True)
    \item $n=4$: $15$ not prime $\Leftrightarrow 4$ not prime. (True)
    \item $n=5$: $31$ is prime $\Leftrightarrow 5$ is prime. (True)
    \item $n=6$: $63$ not prime $\Leftrightarrow 6$ not prime. (True)
    \item $n=11$: $2047$ ($23 \times 89$) not prime $\Leftrightarrow 11$ is prime. (False)
\end{itemize}
The biconditional is true for $n \in$ \textbf{$\{2, 3, 4, 5, 6\}$}.

\newpage
\section{Part 3: In-Depth Analysis of Problems and Techniques}

\subsection{A) Problem Types and General Approach}
The homework problems generally fall into three distinct types:
\begin{itemize}
    \item \textbf{Type 1: Compound Truth Value Evaluation (e.g., 2.21, 2.22).} 
    \textit{Approach:} Determine the standalone truth values of the atomic statements $P, Q, R$. Then apply the strict Boolean logic truth tables step-by-step (innermost parentheses first) to deduce the truth of the larger compound statement.
    \item \textbf{Type 2: Open Sentences over a Discrete Domain (e.g., 2.39, 2.43, 2.44).} 
    \textit{Approach:} For finite domains, substitute each element $x \in S$ into the expressions $P(x)$ and $Q(x)$ and map the Boolean output. This relies heavily on algebraic verification (e.g. factoring polynomials, evaluating exponentials).
    \item \textbf{Type 3: Set Theory and Abstract Deductions (e.g., 2.23, 2.24, 2.27).} 
    \textit{Approach:} Use logic to traverse relationships between abstract sets. For logical deduction puzzles (like 2.27), hypothesize one statement being true and follow the chain of implications until a contradiction arises, ruling out invalid states.
\end{itemize}

\subsection{B) Key Algebraic and Logical Manipulations}
\begin{itemize}
    \item \textbf{Vacuous Truth Enforcement}: Applied heavily (e.g., Ex 2.32a). Recognizing that if an open sentence evaluates to False at $x$, any implication branching from it is automatically True.
    \item \textbf{Algebraic Equivalence Translation}: Converting absolute value statements to set bounds (Ex 2.37). $|x-3| < 1 \Rightarrow -1 < x - 3 < 1 \Rightarrow 2 < x < 4$. 
    \item \textbf{Counter-Example Generation}: Used in evaluating implications. Proving $P \Rightarrow Q$ over a set requires showing $P(x) \Rightarrow Q(x)$ is True for all $x$. A single state where $P(x)$ is True and $Q(x)$ is False refutes it (Ex 2.31a).
\end{itemize}

\section{Part 4: "Cheatsheet" and Tips for Success}

\subsection{Essential Formulas and Equivalences}
\begin{itemize}
    \item \textbf{Implication Definition}: $P \Rightarrow Q \equiv (\sim P) \vee Q$
    \item \textbf{Contrapositive}: $P \Rightarrow Q \equiv (\sim Q) \Rightarrow (\sim P)$
    \item \textbf{Biconditional}: $P \Leftrightarrow Q \equiv (P \Rightarrow Q) \wedge (Q \Rightarrow P)$
    \item \textbf{Negation of Implication}: $\sim (P \Rightarrow Q) \equiv P \wedge (\sim Q)$
\end{itemize}

\subsection{Tips, Tricks, and Pitfalls}
\begin{itemize}
    \item \textbf{The Vacuous Truth Trap}: The most common mistake in logic. When a hypothesis is False, the implication is TRUE. ``If $0=1$, then I am a billionaire'' is a logically true statement.
    \item \textbf{Sufficient vs Necessary}: ``$P$ is sufficient for $Q$'' means $P \Rightarrow Q$. ``$P$ is necessary for $Q$'' means $Q \Rightarrow P$.
    \item \textbf{Converse is NOT Equivalent}: Do not assume $P \Rightarrow Q$ guarantees $Q \Rightarrow P$. That requires a biconditional.
    \item \textbf{Subset Testing Cheat}: For open sentences over a domain, $P(x) \Rightarrow Q(x)$ is universally true if and only if the truth set of $P$ is a \textit{subset} of the truth set of $Q$.
\end{itemize}

\section{Part 5: Conceptual Synthesis and The "Big Picture"}

\subsection{A) Thematic Connections}
The core theme of logic is \textbf{determinism and state definition}, which overlaps heavily with Set Theory and Boolean Algebra. The logical ``AND'' ($\wedge$) acts algebraically like multiplication or set intersection ($\cap$). The logical ``OR'' ($\vee$) acts like addition or set union ($\cup$). The implication $P \Rightarrow Q$ corresponds perfectly to the set notation $P \subseteq Q$.

\subsection{B) Forward and Backward Links}
Backward, this extends the fundamental algebra of truth tables. Forward, this establishes the bedrock for all upper-level math. Understanding implications and contrapositives is mandatory for constructing \textbf{mathematical proofs}. For instance, in Real Analysis, proving limits (the epsilon-delta definition) relies entirely on a complex chain of nested implications.

\newpage
\section{Part 6: Real-World Application and Modeling}

\subsection{A) Concrete Scenarios in Financial Engineering}
\begin{itemize}
    \item \textbf{Derivative Pricing (Options Exercising)}: Logic heavily governs the stochastic modeling of American Options. An algorithm must continuously evaluate the conditional statement: ``If the current intrinsic value of the option is greater than the continuation value determined by the Longstaff-Schwartz model, then the option is exercised early.'' This is a rigid $P \Rightarrow Q$ logical trigger that prices the security.
    \item \textbf{Algorithmic Trading Systems}: Quants build complex logic gates. ``If the 50-day moving average crosses the 200-day moving average AND the RSI is below 30, then execute a buy order.'' $(M \wedge R) \Rightarrow B$.
    \item \textbf{Credit Default Swaps (CDS)}: A CDS triggers a payout if and only if a specific set of default criteria are met. This requires strict biconditional ($P \Leftrightarrow Q$) legal and mathematical modeling.
\end{itemize}

\subsection{B) Model Problem Setup: Credit Default Swap Payout}
Let us model the conditions under which a financial institution must execute a CDS payout.
\begin{itemize}
    \item \textbf{Domain}: $t \in [0, T]$, where $T$ is the maturity date of the contract.
    \item \textbf{Variables}:
    \begin{itemize}
        \item $D(t)$: The underlying corporate bond formally defaults at time $t$.
        \item $G(t)$: The legal 30-day grace period has expired by time $t$.
        \item $E(t)$: The total exposure exceeds the contract threshold at time $t$.
        \item $P(t)$: A payout is legally triggered at time $t$.
    \end{itemize}
\end{itemize}
\textbf{The Mathematical Model:}
To find the state of the payout, the algorithm evaluates the following logical equivalence for all $t$:
\[ P(t) \Leftrightarrow \Big( D(t) \wedge G(t) \wedge E(t) \Big) \]
If the contract required only default OR extreme exposure, the model would shift to:
\[ P(t) \Leftrightarrow \Big( (D(t) \wedge G(t)) \vee E(t) \Big) \]

\section{Part 7: Common Variations and Untested Concepts}

The homework set provided did not thoroughly test two crucial logical equivalents often found in standard curricula:
\begin{itemize}
    \item \textbf{Logical Equivalence Proofs without Truth Tables}: Proving two statements are equivalent using De Morgan's Laws and the definition of implication.
    \item \textbf{The Contrapositive}: Testing the equivalence $P \Rightarrow Q \equiv \sim Q \Rightarrow \sim P$.
\end{itemize}

\textbf{Worked Example (Missing Concept): Prove $P \Rightarrow (Q \vee R)$ is equivalent to $(P \wedge \sim Q) \Rightarrow R$.} \\
\textit{Solution Using Logical Laws}:
\begin{align*}
    P \Rightarrow (Q \vee R) &\equiv \sim P \vee (Q \vee R) \quad &&\text{(Definition of Implication)} \\
    &\equiv (\sim P \vee Q) \vee R \quad &&\text{(Associativity of } \vee) \\
    &\equiv \sim (P \wedge \sim Q) \vee R \quad &&\text{(De Morgan's Law)} \\
    &\equiv (P \wedge \sim Q) \Rightarrow R \quad &&\text{(Definition of Implication)}
\end{align*}
This confirms they are logically identical without needing an 8-row truth table!

\newpage
\section{Part 8: Advanced Diagnostic Testing: "Find the Flaw"}

Below are 5 advanced logic problems with step-by-step solutions that contain a subtle, critical flaw. You must act as the grader, identify the error, and provide the correct steps.

\vspace{10pt}
\textbf{Problem 1:} Let $P(x): x^2 < 0$ and $Q(x): 3x = 5$ over the domain $\mathbb{R}$. Evaluate the truth value of the implication $P(x) \Rightarrow Q(x)$ for all $x \in \mathbb{R}$.
\begin{itemize}
    \item \textbf{Flawed Solution}: Because the square of a real number is never negative, $P(x)$ is always False. Since the hypothesis is False, the implication fails and is False for all $x$.
    \item \textbf{Error}: The student confused a false hypothesis with a false implication.
    \item \textbf{Correction}: In formal logic, a False hypothesis makes the conditional statement vacuously True. Thus, $P(x) \Rightarrow Q(x)$ is \textbf{True} for all $x \in \mathbb{R}$.
\end{itemize}

\vspace{10pt}
\textbf{Problem 2:} Find the negation of the statement: ``If it is raining, then the game is canceled.'' ($P \Rightarrow Q$).
\begin{itemize}
    \item \textbf{Flawed Solution}: The negation of an implication is formed by negating both parts. So the negation is ``If it is not raining, then the game is not canceled'' ($\sim P \Rightarrow \sim Q$).
    \item \textbf{Error}: The student provided the \textit{inverse} of the statement, not the logical negation.
    \item \textbf{Correction}: The negation of $P \Rightarrow Q$ is $P \wedge \sim Q$. The correct negation is: ``It is raining, and the game is not canceled.''
\end{itemize}

\vspace{10pt}
\textbf{Problem 3:} Determine the truth value of the biconditional $P \Leftrightarrow Q$ given that $P$ is True and $Q$ is False.
\begin{itemize}
    \item \textbf{Flawed Solution}: A biconditional means either $P \Rightarrow Q$ or $Q \Rightarrow P$ is True. Since $Q \Rightarrow P$ is (False $\Rightarrow$ True), which evaluates to True, the entire biconditional is True.
    \item \textbf{Error}: The student evaluated the biconditional as an ``OR'' ($\vee$) instead of an ``AND'' ($\wedge$).
    \item \textbf{Correction}: $P \Leftrightarrow Q \equiv (P \Rightarrow Q) \wedge (Q \Rightarrow P)$. While $Q \Rightarrow P$ is True, $P \Rightarrow Q$ is False. True $\wedge$ False yields \textbf{False}.
\end{itemize}

\vspace{10pt}
\textbf{Problem 4:} Find the truth set for the open sentence $x = 3 \Rightarrow |x| = 3$ over $\mathbb{R}$.
\begin{itemize}
    \item \textbf{Flawed Solution}: If $x = 3$, then $|3| = 3$, which is True. Therefore, the truth set is $\{3\}$.
    \item \textbf{Error}: The student only evaluated the case where the hypothesis was True, ignoring the rest of the domain.
    \item \textbf{Correction}: For any $x \neq 3$, the hypothesis $x = 3$ is False, making the implication vacuously True. Thus, the statement is True for all real numbers. The truth set is \textbf{$\mathbb{R}$}.
\end{itemize}

\vspace{10pt}
\textbf{Problem 5:} Evaluate whether the following compound statement is a tautology: $(P \wedge Q) \Rightarrow P$.
\begin{itemize}
    \item \textbf{Flawed Solution}: If $P$ is False and $Q$ is True, then $(P \wedge Q)$ is False. The implication becomes False $\Rightarrow$ False, which is False. Since it can be False, it is not a tautology.
    \item \textbf{Error}: The student incorrectly evaluated False $\Rightarrow$ False. 
    \item \textbf{Correction}: False $\Rightarrow$ False evaluates to \textbf{True}. Constructing the truth table shows that in all four cases, the statement evaluates to True, making it a \textbf{tautology}.
\end{itemize}

\end{document}